% =====================================================================
% Appendix B: CST-Bench-Syn construction and data card (one-column appendix).
% PLANNED: nothing of CST-Bench-Syn is built yet. The text specifies the design;
% tab:syn has blank cells for measured values and marks plan targets as targets.
% Never state plan estimates (hours, voice counts, cost) as facts.
% Holds the Syn detail that Sec. 4.1 leaves out: sources and their
% translation provenance, Korean consensus MT, TTS editions and QA,
% calibration split, translationese rule, release and licences, and the
% voice-consent rule (decision: "clone no real person's voice without
% consent"; Sec. 4 no longer repeats it, the Impact Statement does).
% =====================================================================
\section{CST-Bench-Syn Construction and Data Card}\label{app:syn}

\draftnote{planned: \syn is not built yet; this appendix specifies its construction, and every measured value in \cref{tab:syn} is pending}

\syn is an evaluation-only set of synthetic (TTS) speech. It adds what \taxi cannot provide: turn-end overlap that can be measured at the word level, a distant language pair, and both speaker--language assignments, which remove TAXI's first-speaker prior and its tie between language and role. It is not real conversation, does not meet the requirements for a real test set (\cref{app:taxi}) and is translationese on one side of every pair. We therefore draw only relative conclusions from it (gaps between protocols and system rankings, never absolute scores across directions or pairs) and test whether it ranks systems as \taxi does (\cref{tab:validity,fig:sim2real}). There is no training or development split: a calibration split of about 5\% of the sessions, disjoint from the evaluation sessions, is the only data on which wrapper thresholds and operating points are chosen, with one rule for all systems. \Cref{tab:syn} is the data card.

% =====================================================================
% Table A2 (tab:syn): CST-Bench-Syn data card. Included from App. B.
% PLANNED: the set is not built. Measured cells use \tbd (blank); the last
% column holds plan targets (not results). Known now: reference type, licences.
% Do not add hours, voice counts or costs: they are plan estimates, not facts.
% =====================================================================
\begin{table}[htbp]
  \caption{\syn data card (planned). The set is not built yet: blank cells will hold measured values, and the last column gives plan targets, which are not results. A session is one dialogue under one speaker--language assignment; every dialogue is rendered under both. The re-recognition error is the WER for English and German and the CER for Korean; turns above the target are re-synthesised. MT: machine translation; n/a: not applicable; --: no target.}
  \label{tab:syn}
  \centering
  \footnotesize
  \setlength{\tabcolsep}{4pt}
  \begin{tabular}{@{}lccccc@{}}
    \toprule
    & \multicolumn{2}{c}{\pair{En}{De}} & \multicolumn{2}{c}{\pair{Ko}{En}} & \\
    \cmidrule(lr){2-3}\cmidrule(lr){4-5}
    Statistic & XDailyDialog & BConTrasT & XDailyDialog & BConTrasT & Plan target \\
    \midrule
    \multicolumn{6}{@{}l}{\textit{Size}} \\
    Dialogues\,/\,sessions                    & \tbd & \tbd & \tbd & \tbd & -- \\
    Turns per direction                       & \tbd & \tbd & \tbd & \tbd & -- \\
    Speech per timing condition (h)           & \tbd & \tbd & \tbd & \tbd & -- \\
    Calibration split (\% of sessions)        & \tbd & \tbd & \tbd & \tbd & about 5 \\
    \midrule
    \multicolumn{6}{@{}l}{\textit{Speech}} \\
    Voices per language and edition           & \tbd & \tbd & \tbd & \tbd & -- \\
    Re-recognition error (\%)                 & \tbd & \tbd & \tbd & \tbd & $\leq 5$ \\
    Re-synthesised turns (\%)                 & \tbd & \tbd & \tbd & \tbd & -- \\
    Speaker cosine, same\,/\,other speaker    & \tbd & \tbd & \tbd & \tbd & -- \\
    TTS word times vs.\ forced alignment      & \tbd & \tbd & \tbd & \tbd & -- \\
    \midrule
    \multicolumn{6}{@{}l}{\textit{Text}} \\
    Reference type                            & human & human & \makecell{Ko: consensus MT\\En: human} & \makecell{Ko: consensus MT\\En: human} & -- \\
    Korean consensus pass rate (\%)           & n/a  & n/a  & \tbd & \tbd & $\geq 85$ \\
    Korean spot-check error rate (\%)         & n/a  & n/a  & \tbd & \tbd & -- \\
    Contamination-probe agreement             & \tbd & \tbd & \tbd & \tbd & -- \\
    Licence                                   & CC BY-NC-SA 4.0 & CC BY-SA 4.0 & CC BY-NC-SA 4.0 & CC BY-SA 4.0 & -- \\
    \bottomrule
  \end{tabular}
\end{table}


\paragraph{Pairs and speaker assignment.}
\pair{En}{De} is the main pair: it shares its languages with \taxi, which makes the sim-to-real comparison possible. \pair{Ko}{En} adds a distant word order. Every dialogue is rendered under both speaker--language assignments (speaker $A$ speaks the first language and $B$ the second, and vice versa), and the same dialogues serve both pairs, so differences between pairs come from the languages rather than from the content. We exclude German--Korean, which almost no baseline supports, and use English as the pivot.

\paragraph{Dialogues.}
XDailyDialog \citep{liu2023xdailydialog} provides professional translations of the everyday dialogues of DailyDialog \citep{li2017dailydialog}, and BConTrasT \citep{farajian2020findings} provides English--German versions of task-oriented Taskmaster-1 dialogues \citep{byrne2019taskmaster} between a customer and an agent in six domains, machine-translated into German and post-edited by native German speakers. Speaker $A$'s turns come from one language version of a dialogue and $B$'s from the other, so both share one context and the parallel version of each turn is its reference. We select dialogues of 6--16 turns and balance the domains, the turn lengths and the share of dialogues that contain numbers or named entities. XDailyDialog turns are short (about 11.5 tokens on average), so the selection favours multi-sentence turns, which leave room for turn-end overlap. No parallel corpus that cannot be redistributed is used.

\paragraph{Contamination.}
DailyDialog and Taskmaster-1 text may be in the pretraining data of the LLMs that serve as baselines and teachers. For each of them, a probe gives the opening turns of a selected dialogue, lets the model continue it and measures the agreement of the continuation with the original (\cref{tab:syn}). Generation prompts, models and seeds are released. No DailyDialog, Taskmaster-1 or TAXI text and no \syn dialogue or voice is used to train \model (\cref{sec:model}).

\paragraph{References.}
English--German references are the parallel versions of each turn: the English side is original, and the German side is a professional translation (XDailyDialog) or a post-edited machine translation (BConTrasT). Neither source contains Korean, so every Korean text, both the Korean turns that are synthesised and the Korean references of English turns, is translated from English by three independent systems: Qwen3.8-27B, MADLAD-400-10B-MT \citep{kudugunta2023madlad} and EXAONE-4.0 \citep{lgai2025exaone4}. A text passes if at least two of the three translations agree, as judged by MetricX-24 in its reference-free mode \citep{juraska2024metricx}, CometKiwi \citep{rei2022cometkiwi}, the LaBSE similarity \citep{feng2022labse} of back-translations, and matching numbers and named entities; text that fails is removed, and we report the pass rate. The released Korean text is the MADLAD-400 translation, as no evaluated system belongs to that model family: Qwen3.8-27B is also a baseline translator and the teacher of our commit targets (\cref{sec:model}), and EXAONE-4.0, whose licence restricts its outputs to non-commercial use, serves only in the agreement check, so its outputs are neither released nor used for training.\draftnote{fix the rule for text whose MADLAD-400 translation is the dissenting one}
Native Korean speakers annotate the errors in a fixed random sample with a light MQM-style scheme \citep{freitag2021experts}. \pair{Ko}{En} enters the main results only if at least 85\% of the Korean text passes and the speech meets the re-recognition target below; otherwise it is reported as a pilot. For \pair{En}{De}, silver references built in the same way test whether gold and silver references rank systems alike (\cref{tab:validity}). Gold and silver segments are aggregated separately, and reference-free CometKiwi is reported next to the reference-based metrics.

\paragraph{Speech synthesis and quality assurance.}
A commercial edition carries the main results: Azure neural TTS \citep{microsoft2026text} with native preset voices for en-US, de-DE and ko-KR, balanced by gender, and word boundaries from the synthesiser. An open edition, CosyVoice~3 \citep{du2025cosyvoice3} or Qwen3-TTS \citep{hu2026qwen3tts} (both with Apache-2.0 weights), covers a subset and tests whether the TTS engine changes system rankings (\cref{tab:validity}). A pilot of 50 sentences per language admits an engine only if its re-recognition error is at most 5\%. Each speaker keeps one voice throughout a dialogue. Voices are presets, designed voices or voices cloned from CC0 reference recordings; we clone no real person's voice without consent and no voice across languages. Every synthesised turn is re-recognised with Whisper-large-v3 \citep{radford2023robust} and Qwen3-ASR \citep{shi2026qwen3asr}; turns with a WER (CER for Korean) above 5\% are re-synthesised, and we report the error and the re-synthesis rate. Speaker consistency is measured by the cosine similarity of speaker embeddings within and between speakers. The synthesiser's word boundaries are compared with forced alignment (Qwen3-ForcedAligner; \citealp{shi2026qwen3asr}), and we report the distribution of their differences, because word-level overlap and the reference points of latency depend on these times.

\paragraph{Rendering and release.}
Sessions are rendered with the same renderer and rules as \taxi (\cref{sec:bench}) under L0 natural and L1, optionally also under L0 mediated, each with its own configuration name and checksums; overlap is measured on turn intervals and on word boundaries. Mixes are clean: noise and reverberation variants are outside this release, and a telephone-band variant is added only if \syn and \taxi rank systems differently. Each session comes with the mono mix, per-speaker channels (for analysis and oracle diagnostics only), a timeline with turn and word times, and references flagged as gold or silver. A data card discloses that the speech is synthetic and lists the sources, licences, generation models, prompts, seeds and QA results. CC BY-NC-SA 4.0 and CC BY-SA 4.0 are incompatible ShareAlike licences, so the XDailyDialog-based and the BConTrasT-based sessions form two subsets, each under the licence of its source (Taskmaster-1, on which BConTrasT builds, is CC BY 4.0). The commercial edition is synthesised only on a paid tier (\cref{tab:licences}). Instead of a fixed size in hours, \syn is sized by a power estimate from the variance of TAXI segment scores for the main comparisons, fixed before the test runs: a first tranche is built and extended only if its confidence intervals are too wide.

\paragraph{Out of scope.}
This release leaves out the parts of an earlier design that targeted interruptions and barge-in: episodes of two to four history turns, an interruption event and one or two resolution turns; information-causal event generation with three difficulty levels; sizing by at least 300 events per combination of event type, difficulty and pair, which implied about 10--12\,h per pair; partial references for interrupted turns; leakage checks for generated events; and a high-overlap condition with at least 25\% overlap. Interruptions remain future work (\cref{sec:conclusion}).

\paragraph{Limitations.}
TTS prosody lacks disfluencies, reactions to the partner and the Lombard effect; the dialogue text is largely planned rather than spontaneous; one side of every pair is translationese; the Korean references are silver; and wrapper thresholds calibrated on clean 16\,kHz TTS may suit telephone-band \taxi less well. Rank agreement with \taxi and between the two TTS editions (\cref{tab:validity}) indicates how far conclusions drawn on \syn transfer to real speech.
